\renewcommand{\bibname}{Sources and reading notes}
\begin{thebibliography}{9}
\addcontentsline{toc}{chapter}{Sources and reading notes}
\bibitem{iupac-conformation}
IUPAC--IUB Joint Commission on Biochemical Nomenclature.
\emph{Abbreviations and Symbols for the Description of Conformations of
Polynucleotide Chains}, Recommendations 1982, published 1983.
\url{https://iupac.qmul.ac.uk/misc/pnuc1.html}.
Relevant nomenclature body inspected: atom numbering, chain direction,
and internucleotide connectivity. Our drawings are simplified topological
projections, not reproductions of the source figures.

\bibitem{iupac-symbols}
IUPAC--IUB. \emph{Abbreviations and Symbols for Nucleic Acids,
Polynucleotides and their Constituents}, Recommendations 1970.
\url{https://iupac.qmul.ac.uk/misc/naabb.html}.
Relevant nucleoside/nucleotide and phosphate-position conventions inspected.
Modern ambiguity symbols are taken from INSDC, not inferred from this
historical nomenclature document.

\bibitem{nhgri}
National Human Genome Research Institute.
\emph{Deoxyribonucleic Acid (DNA) Fact Sheet}, updated 24 August 2020.
\url{https://www.genome.gov/about-genomics/fact-sheets/Deoxyribonucleic-Acid-Fact-Sheet}.
Full fact-sheet body inspected. Used for basic components, complementary
pairing, and opposing strand directions; not as a quantitative binding model.

\bibitem{insdc}
International Nucleotide Sequence Database Collaboration.
\emph{Feature Table Definition}, Appendix 7.4.1, nucleotide base codes.
\url{https://www.insdc.org/submitting-standards/feature-table/}.
The complete relevant base-code appendix was inspected. The generated table
and set-complement algorithm are our own implementation of those code meanings.

\bibitem{stacking}
P. Yakovchuk, E. Protozanova, and M. D. Frank-Kamenetskii (2006).
Base-stacking and base-pairing contributions into thermal stability of the
DNA double helix. \emph{Nucleic Acids Research} 34(2), 564--574.
\url{https://doi.org/10.1093/nar/gkj454}.
Access was limited to indexed article text; direct PMC access encountered
a challenge page. Used only for the qualitative distinction between
stacking and pairing, not numerical parameter values.

\bibitem{stacking-correction}
Correction notice for the preceding article (2006).
\emph{Nucleic Acids Research} 34(3), 1082.
\url{https://doi.org/10.1093/nar/gkj523}.
The indexed notice reports errors in Figures 5--8 and republication.
We did not independently revalidate the corrected figures and reproduce
none of their curves, values, or decompositions.
\end{thebibliography}
